% Part:sets-functions-relations
% Chapter: size-of-sets
% Section: non-enumerability

\documentclass[../../../include/open-logic-section]{subfiles}

\begin{document}

\olfileid[fa-IR]{sfr}{siz}{nen}

\olsection{مجموعه‌های \printtoken{S}{nonenumerable}}

\begin{editorial}
  این بخش ناشمارابودنِ $\Bin^\omega$ و $\Pow{\PosInt}$ را با استفاده
  از تعریفِ \olref[enm]{sec} ثابت می‌کند. این بخش طوری
  تنظیم شده است که اندکی مقدماتی‌تر و اندکی مشروح‌تر از روایتِ
  \olref[enm-alt]{sec} باشد
\end{editorial}

برخی مجموعه‌ها، مانند مجموعهٔ~$\PosInt$ از اعداد صحیح مثبت،
نامتناهی‌اند. تاکنون نمونه‌هایی از مجموعه‌های نامتناهی دیده‌ایم که همگی
!!{enumerable} بوده‌اند. بااین‌حال، مجموعه‌های نامتناهی‌ای نیز وجود دارند
که این ویژگی را ندارند. چنین مجموعه‌هایی را \emph{!!{nonenumerable}}
می‌نامند.

پیش از هر چیز، شاید همین وجود مجموعه‌های !!{nonenumerable} شگفت‌آور
باشد. برای هر مجموعهٔ !!{enumerable} مانند~$A$ تابعی !!a{surjective}
به‌صورت $f \colon \PosInt \to A$ وجود دارد. اگر مجموعه‌ای
!!{nonenumerable} باشد، چنین تابعی وجود ندارد. یعنی هیچ تابعی که
!!{element}s را از مجموعهٔ نامتناهیِ~$\PosInt$ به~$A$ نگاشت کند،
نمی‌تواند همهٔ~$A$ را فراگیرد. پس شمارِ !!{element}s در~$A$ از شمار
بی‌نهایت عدد صحیح مثبت نیز «بیشتر» است.

چگونه می‌توان ثابت کرد که مجموعه‌ای !!{nonenumerable} است؟ باید نشان
دهید که هیچ تابع پوشایی از این نوع نمی‌تواند وجود داشته باشد. به‌طور
هم‌ارز، باید نشان دهید که نمی‌توان عناصر~$A$ را در فهرستی نامتناهی و
یک‌سویه برشمرد. بهترین راه برای این کار آن است که نشان دهید هر فهرستی
از !!{element}sی از~$A$ ناگزیر دست‌کم یک عضو را از قلم می‌اندازد؛ یا
اینکه هیچ تابع $f\colon \PosInt \to A$ نمی‌تواند !!{surjective} باشد.
می‌توانیم این کار را با \emph{روش قطری} کانتور انجام دهیم. با داشتن
فهرستی از !!{element}sی از~$A$، مثلاً $x_1$، $x_2$، \dots، عضو دیگری
از~$A$ می‌سازیم که بنا بر نحوهٔ ساختش، به‌هیچ‌وجه نمی‌تواند در آن
فهرست باشد.

نخستین مثال ما مجموعهٔ~$\Bin^\omega$ از همهٔ دنباله‌های نامتناهی و
بدون شکافِ $0$ها و $1$هاست.

\begin{thm}
\ollabel{thm:nonenum-bin-omega}
$\Bin^\omega$~!!{nonenumerable} است.
\end{thm}

\begin{proof}
به قصد خلف، فرض کنید $\Bin^\omega$ !!{enumerable} است؛ یعنی فرض کنید
$s_{1}$، $s_{2}$، $s_{3}$، $s_{4}$، \dots{} فهرستی است که همهٔ
!!{element}s را از~$\Bin^\omega$ در بر می‌گیرد. هر $s_i$ خود دنباله‌ای
نامتناهی از $0$ها و~$1$هاست. عضو $j$اُمِ دنبالهٔ $i$اُم در این فهرست
را $s_i(j)$ می‌نامیم. پس دنبالهٔ $i$اُم، یعنی~$s_i$، چنین است:
\[
s_i(1), s_i(2), s_i(3), \dots
\]

می‌توانیم این فهرست و اعضای هر دنبالهٔ $s_i$ در آن را در آرایه‌ای
مرتب کنیم:
\[
\begin{array}{c|c|c|c|c|c}
& 1 & 2 & 3 & 4 & \dots \\\hline
1 & \mathbf{s_{1}(1)} & s_{1}(2) & s_{1}(3) & s_1(4) & \dots \\\hline
2 & s_{2}(1)& \mathbf{s_{2}(2)} & s_2(3) & s_2(4) & \dots \\\hline
3 & s_{3}(1)& s_{3}(2) & \mathbf{s_3(3)} & s_3(4) & \dots \\\hline
4 & s_{4}(1)& s_{4}(2) & s_4(3) & \mathbf{s_4(4)} & \dots \\\hline
\vdots & \vdots & \vdots & \vdots & \vdots & \mathbf{\ddots}
\end{array}
\]
برچسب‌های کنار آرایه شمارهٔ دنباله‌ها را در فهرستِ $s_1$، $s_2$،
\dots مشخص می‌کنند؛ اعداد بالای آرایه، !!{element}sی را که در هریک از
دنباله‌ها قرار دارند مشخص می‌کنند. برای نمونه، $s_{1}(1)$ نامِ عددی،
یعنی یک $0$ یا یک~$1$، است که نخستین !!{element} در دنبالهٔ $s_{1}$
است؛ و همین‌طور تا آخر.

اکنون دنبالهٔ نامتناهیِ $\overline{s}$ را از $0$ها و $1$ها می‌سازیم؛
دنباله‌ای که به‌هیچ‌وجه نمی‌تواند در این فهرست باشد. تعریف
$\overline{s}$ به فهرستِ $s_1$، $s_2$، \dots وابسته خواهد بود. هر
فهرست نامتناهی از دنباله‌های نامتناهیِ $0$ها و $1$ها، دنباله‌ای
نامتناهی مانند~$\overline{s}$ پدید می‌آورد که تضمین شده است در آن
فهرست ظاهر نشود.

برای تعریف $\overline{s}$، تمام !!{element}s را یک‌به‌یک مشخص می‌کنیم؛
یعنی $\overline{s}(n)$ را برای همهٔ $n \in \PosInt$ تعیین می‌کنیم.
برای این کار قطر آرایهٔ بالا را از بالا به پایین می‌خوانیم (نام «روش
قطری» نیز از همین‌جا می‌آید) و سپس هر $1$ را به $0$ و هر $0$ را به~$1$
تغییر می‌دهیم. به‌بیان انتزاعی‌تر، $\overline{s}(n)$ را برابر $0$ یا
$1$ تعریف می‌کنیم، بسته به اینکه !!{element} $n$اُمِ قطر، یعنی
$s_n(n)$، برابر $1$ یا $0$ باشد.
\[
\overline{s}(n) =
\begin{cases}
1 & \text{اگر $s_{n}(n) = 0$}\\
0 & \text{اگر $s_{n}(n) = 1$}.
\end{cases}
\]
اگر فرمول‌ها را بر تعریف موردی ترجیح می‌دهید، می‌توانید
$\overline{s}(n) = 1 - s_n(n)$ را نیز به‌عنوان تعریف به کار ببرید.

آشکارا $\overline{s}$ دنباله‌ای نامتناهی از $0$ها و $1$هاست، زیرا
صرفاً دنبالهٔ مکملِ دنباله‌ای از $0$ها و $1$هاست که بر قطر آرایهٔ ما
ظاهر می‌شوند. پس $\overline{s}$ !!a{element} از~$\Bin^\omega$ است.
اما نمی‌تواند در فهرستِ $s_1$، $s_2$، \dots{} باشد. چرا؟

این دنباله نمی‌تواند نخستین دنبالهٔ فهرست، یعنی $s_1$، باشد، زیرا در
نخستین !!{element} با $s_1$ تفاوت دارد. مقدار $s_1(1)$ هرچه باشد،
$\overline{s}(1)$ را مخالف آن تعریف کرده‌ایم. همچنین نمی‌تواند دومین
دنبالهٔ فهرست باشد، زیرا $\overline{s}$ در عضو دوم با $s_2$ تفاوت
دارد: اگر $s_2(2)$ برابر $0$ باشد، $\overline{s}(2)$ برابر $1$ است و
برعکس. همین استدلال ادامه می‌یابد.

دقیق‌تر بگوییم: اگر $\overline{s}$ در فهرست بود، عددی مانند $k$ وجود
داشت که $\overline{s} = s_{k}$. دو دنباله یکسان‌اند اگر و تنها اگر در
هر جایگاه برابر باشند؛ یعنی برای هر~$n$، $\overline{s}(n) =
s_{k}(n)$. پس به‌ویژه، با گرفتن حالت خاصِ $n = k$، باید
$\overline{s}(k) = s_{k}(k)$ برقرار باشد. $s_k(k)$ یا $0$ است یا~$1$.
اگر $0$ باشد، $\overline{s}(k)$ باید~$1$ باشد---زیرا $\overline{s}$
را چنین تعریف کردیم. اما اگر $s_k(k) = 1$ باشد، باز هم به‌سبب نحوهٔ
تعریف $\overline{s}$، داریم $\overline{s}(k) = 0$. در هر دو حالت
$\overline{s}(k) \neq s_{k}(k)$.

کار را با این فرض آغاز کردیم که فهرستی از !!{element}s در
$\Bin^\omega$، یعنی $s_1$، $s_2$، \dots{}، وجود دارد. از این فهرست
دنبالهٔ~$\overline{s}$ را ساختیم و ثابت کردیم که نمی‌تواند در فهرست
باشد. بااین‌حال، این نیز بی‌تردید دنباله‌ای از $0$ها و $1$هاست، اگر
همهٔ $s_i$ها دنباله‌هایی از $0$ها و $1$ها باشند؛ یعنی $\overline{s} \in
\Bin^\omega$. این نتیجه به‌ویژه نشان می‌دهد که هیچ
فهرستی از \emph{همهٔ} !!{element}s در~$\Bin^\omega$ نمی‌تواند وجود داشته
باشد، زیرا از هر فهرست این‌چنینی می‌توانستیم دنبالهٔ~$\overline{s}$ را
بسازیم که تضمین شده است در فهرست نباشد. بنابراین فرض وجود فهرستی از
همهٔ دنباله‌های~$\Bin^\omega$ به تناقض می‌انجامد.
\end{proof}

\begin{explain}
این روش برهان را «قطری‌سازی» می‌نامند، زیرا برای تعریف~$\overline{s}$
از قطر آرایه استفاده می‌کند. قطری‌سازی لزوماً به وجود آرایه نیاز
ندارد: می‌توانیم با اندیشه‌ای مشابه نشان دهیم که مجموعه‌ها
!!{enumerable} نیستند، حتی اگر هیچ آرایه و هیچ قطر واقعی‌ای در میان
نباشد.
\end{explain}

\begin{thm}
\ollabel{thm:nonenum-pownat}
$\Pow{\PosInt}$ !!{enumerable} نیست.
\end{thm}

\begin{proof}
به همان شیوه پیش می‌رویم: نشان می‌دهیم که برای هر فهرستی از زیرمجموعه‌های
~$\PosInt$، زیرمجموعه‌ای از $\PosInt$ وجود دارد که نمی‌تواند در فهرست
باشد. فرض کنید فهرست زیر از زیرمجموعه‌های~$\PosInt$ داده شده است:
\[
Z_{1}, Z_{2}, Z_{3}, \dots
\]
اکنون مجموعه‌ای مانند $\overline{Z}$ چنان تعریف می‌کنیم که برای هر
$n \in \PosInt$، $n \in \overline{Z}$ اگر و تنها اگر $n \notin Z_{n}$:
\[
\overline{Z} = \Setabs{n \in \PosInt}{n \notin Z_n}
\]
$\overline{Z}$ آشکارا مجموعه‌ای از اعداد صحیح مثبت است، زیرا بنا بر
فرض هر~$Z_n$ چنین مجموعه‌ای است، و ازاین‌رو $\overline{Z} \in
\Pow{\PosInt}$. اما $\overline{Z}$ نمی‌تواند در فهرست باشد. برای نشان
دادن این مطلب، ثابت می‌کنیم که برای هر $k \in \PosInt$، $\overline{Z} \neq
Z_k$.

پس $k \in \PosInt$ را دلخواه بگیرید. $\overline{Z}$ را چنان تعریف
کرده‌ایم که برای هر $n \in \PosInt$، $n \in \overline{Z}$ اگر و تنها
اگر $n \notin Z_n$. به‌ویژه، با گرفتن $n=k$، داریم $k \in \overline{Z}$
اگر و تنها اگر $k \notin Z_k$. اما این نشان می‌دهد که
$\overline{Z} \neq Z_k$، زیرا $k$ !!a{element} یکی از آن دو است، اما عضو
آن یکی نیست؛ پس $\overline{Z}$ و $Z_k$ از نظر !!{element}s تفاوت دارند.
ازآنجاکه $k$ دلخواه بود، $\overline{Z}$ در فهرستِ $Z_1$، $Z_2$،
\dots نیست.
\end{proof}

\begin{explain}
برهان پیشین از قطر نام نبرد، اما اگر آن را به شکل زیر تصویر کنید،
می‌توانید قطر را در آن ببینید: مجموعه‌های $Z_1$، $Z_2$، \dots را در
آرایه‌ای بنویسید که در آن هر !!{element} مانند~$j \in Z_i$ در ستون
$j$اُم قرار گیرد. فرض کنید چهار مجموعهٔ نخست در آن فهرست به‌ترتیب
$\{1,2,3,\dots\}$، $\{2, 4, 6, \dots\}$، $\{1,2,5\}$ و
$\{3,4,5,\dots\}$ باشند. در این صورت آرایه چنین آغاز می‌شود:
\[
\begin{array}{r@{}rrrrrrr}
  Z_1 = \{ & \mathbf{1}, & 2, & 3, & 4, & 5, & 6, & \dots\}\\
  Z_2 = \{ &  & \mathbf{2}, &  & 4, &  & 6, & \dots\}\\
  Z_3 = \{ & 1, & 2, &  &  & 5\phantom{,} &  & \}\\
  Z_4 = \{ &  &  & 3, & \mathbf{4}, & 5, & 6, & \dots\}\\
  \vdots & & & & & \ddots
\end{array}
\]
آنگاه $\overline{Z}$ مجموعه‌ای است که با پایین‌رفتن روی قطر به دست
می‌آید: هر عددی را که روی قطر ظاهر می‌شود کنار می‌گذاریم و هر $j$ای را
که آرایه در سطر/ستون $j$اُم شکاف دارد، وارد مجموعه می‌کنیم. در نمونهٔ
بالا، $1$ و $2$ را کنار می‌گذاریم، $3$ را وارد می‌کنیم، $4$ را کنار
می‌گذاریم، و همین‌طور ادامه می‌دهیم.
\end{explain}

\begin{prob}
با یک استدلال قطری نشان دهید که $\Pow{\Nat}$ !!{nonenumerable} است.
\end{prob}

\begin{prob}\label{sfr:siz:nen:prob:f-posint}
با یک استدلال قطریِ صریح نشان دهید که مجموعهٔ توابع
$f \colon \PosInt \to \PosInt$ !!{nonenumerable} است. یعنی نشان دهید
اگر $f_1$، $f_2$، \dots فهرستی از توابع باشد و هر $f_i\colon
\PosInt \to \PosInt$ باشد، آنگاه تابعی مانند $\overline{f}\colon \PosInt \to
\PosInt$ وجود دارد که در این فهرست نیست.
\end{prob}

\end{document}
